\documentclass[11pt,a4paper,fontset=fandol]{ctexart}

\usepackage[a4paper,top=2.4cm,bottom=2.6cm,left=2.35cm,right=2.35cm,headheight=18pt,headsep=0.55cm,footskip=26pt]{geometry}
\usepackage{amsmath,amssymb,amsthm}
\usepackage{fancyhdr}
\usepackage{lastpage}
\usepackage{enumitem}
\usepackage{titlesec}
\usepackage{xcolor}
\usepackage{tikz}
\usepackage{hyperref}
\usepackage{bookmark}
\usepackage[most]{tcolorbox}
\usepackage{setspace}
\usepackage{array}
\usetikzlibrary{arrows.meta,positioning}

\hypersetup{
  colorlinks=true,
  linkcolor=blue!50!black,
  urlcolor=blue!50!black,
  citecolor=blue!50!black
}

\setstretch{1.16}
\setlength{\parindent}{2em}
\setlength{\parskip}{0.35em}
\allowdisplaybreaks
\setlist[enumerate]{itemsep=0.32em, topsep=0.32em, leftmargin=2.3em}
\setlist[itemize]{itemsep=0.25em, topsep=0.25em, leftmargin=2em}
\setcounter{tocdepth}{2}

\definecolor{mainblue}{RGB}{36,83,140}
\definecolor{lightblue}{RGB}{241,246,252}
\definecolor{lightgray}{RGB}{246,246,246}
\definecolor{accent}{RGB}{153,76,0}
\definecolor{rulegray}{RGB}{110,118,128}
\definecolor{softgreen}{RGB}{236,247,240}

\titleformat{\section}
  {\Large\bfseries\color{mainblue}}
  {\thesection.}
  {0.45em}
  {}
  [\vspace{0.25em}\color{rulegray}\titlerule]
\titlespacing*{\section}{0pt}{1.3em}{0.9em}
\titleformat{\subsection}{\large\bfseries\color{mainblue}}{\thesubsection}{0.5em}{}
\titlespacing*{\subsection}{0pt}{1.05em}{0.55em}
\titleformat{\subsubsection}{\normalsize\bfseries\color{accent}}{\thesubsubsection}{0.5em}{}

\pagestyle{fancy}
\fancyhf{}
\fancyhead[L]{\small 组合专题课堂讲义}
\fancyhead[C]{\small 图论入门中的组合计数}
\fancyhead[R]{\small 刘文博}
\fancyfoot[C]{\small 第 \thepage 页 / 共 \pageref{LastPage} 页}
\renewcommand{\headrulewidth}{0.55pt}
\renewcommand{\footrulewidth}{0.35pt}

\newtheoremstyle{formaltheorem}
  {0.8em}{0.8em}{\normalfont}{}{\bfseries\color{mainblue}}{．}{0.6em}{}
\newtheoremstyle{formalexample}
  {0.8em}{0.8em}{\normalfont}{}{\bfseries\color{accent}}{．}{0.6em}{}

\theoremstyle{formaltheorem}
\newtheorem{theorem}{定理}[section]
\newtheorem{method}{方法}[section]
\theoremstyle{formalexample}
\newtheorem{example}{例题}[section]
\newtheorem{exercise}{练习}[section]

\newenvironment{solution}{\par\noindent\textbf{解：}}{\hfill$\square$\par}

\newtcolorbox{goalbox}[1]{
  breakable,
  colback=lightblue,
  colframe=mainblue,
  boxrule=0.7pt,
  arc=1mm,
  left=1.2mm,
  right=1.2mm,
  top=1mm,
  bottom=1mm,
  title=\textbf{#1},
  fonttitle=\bfseries
}

\newtcolorbox{notebox}[1]{
  breakable,
  colback=lightgray,
  colframe=black!50,
  boxrule=0.55pt,
  arc=1mm,
  left=1.2mm,
  right=1.2mm,
  top=1mm,
  bottom=1mm,
  title=\textbf{#1},
  fonttitle=\bfseries
}

\newtcolorbox{skillbox}[1]{
  breakable,
  colback=softgreen,
  colframe=green!40!black,
  boxrule=0.65pt,
  arc=1mm,
  left=1.2mm,
  right=1.2mm,
  top=1mm,
  bottom=1mm,
  title=\textbf{#1},
  fonttitle=\bfseries
}

\begin{document}

\thispagestyle{empty}
\begin{titlepage}
\centering
\vspace*{0.9cm}

\begin{tcolorbox}[
  enhanced,
  width=0.92\textwidth,
  colback=white,
  colframe=mainblue,
  boxrule=1pt,
  arc=0mm,
  left=6mm,
  right=6mm,
  top=8mm,
  bottom=8mm
]
\centering

{\fontsize{24}{30}\selectfont\bfseries 组合专题课堂讲义\par}
\vspace{0.7cm}
{\fontsize{17}{22}\selectfont 图论入门中的组合计数\par}

\vspace{1.1cm}
\color{rulegray}\rule{0.72\textwidth}{0.7pt}\par
\vspace{1cm}

\begin{tcolorbox}[colback=lightblue,colframe=mainblue,width=0.82\textwidth,boxrule=1pt]
\centering
{\Large 适合对象：小学高年级至初中低年级数学竞赛学生}\\[0.4em]
{\large 已具备基础计数能力，准备学习“用图来表达关系并进行计数”}
\end{tcolorbox}

\vspace{1.2cm}

\renewcommand{\arraystretch}{1.42}
\begin{tabular}{>{\bfseries}p{3.5cm}p{8cm}}
讲义作者： & 刘文博 \\
专题关键词： & 图、点、边、度数、握手定理、完全图、二分图、网格路径 \\
建议课时： & 2--3 课时 \\
专题目标： & 学会把“关系”翻译成图，并利用图的结构进行计数与证明 \\
编译方式： & XeLaTeX \\
\end{tabular}

\vspace{1.2cm}
\color{rulegray}\rule{0.72\textwidth}{0.7pt}\par
\vspace{0.8cm}

{\normalsize 本讲从“点和边”最基本的图论语言出发，讲握手定理、完全图与二分图的边数、网格图中的路径计数，以及图论里最常见的计数证明方式。\par}

\vspace{0.5cm}
{\large \today\par}
\end{tcolorbox}
\end{titlepage}

\pagenumbering{Roman}
\pagestyle{plain}
\tableofcontents
\newpage
\pagestyle{fancy}
\pagenumbering{arabic}

\section{专题目标与核心问题}

\begin{goalbox}{这一讲要解决什么问题}
很多竞赛题里，人物之间的关系、城市之间的连接、点之间的通路，都可以看成“图”。

一旦把问题翻译成图，就能用“点数、边数、度数、路径数”来进行计数，有时会比直接在文字题面中思考清楚得多。
\end{goalbox}

\begin{goalbox}{学完以后希望达到的能力}
\begin{itemize}
  \item 知道什么是点、边、度数；
  \item 会使用握手定理 $\sum \deg = 2E$；
  \item 会计算完全图、二分图中的边数；
  \item 会把网格路径问题看成图上的路径计数；
  \item 能用图论语言解释一些组合计数结论。
\end{itemize}
\end{goalbox}

\begin{notebox}{这讲和前面专题的关系}
\begin{itemize}
  \item 和排列组合相比，这里更强调“把关系画成图”；
  \item 和递推计数相比，这里也会出现路径递推；
  \item 和抽屉原理、染色问题相比，图论会把许多结构题统一在一张图上观察。
\end{itemize}
\end{notebox}

\begin{center}
\begin{tikzpicture}[
  scale=0.95,
  every node/.style={circle,draw=mainblue,fill=lightblue,inner sep=1.6pt,font=\small},
  edge/.style={draw=black!65,line width=0.9pt}
]
  \node[draw=none,fill=none,circle=false,font=\bfseries\color{mainblue}] at (2.2,2.9) {图就是“对象 + 关系”};
  \node (A) at (0,1.2) {$A$};
  \node (B) at (1.6,2.2) {$B$};
  \node (C) at (3.2,1.5) {$C$};
  \node (D) at (2.7,-0.3) {$D$};
  \node (E) at (0.7,-0.5) {$E$};
  \foreach \u/\v in {A/B,B/C,C/D,D/E,E/A,B/D} {\draw[edge] (\u)--(\v);}
  \node[draw=none,fill=none,circle=false,font=\scriptsize,align=left] at (5.9,1.2) {点：对象\\边：关系\\度数：连出去几条边};
\end{tikzpicture}
\end{center}

\section{什么是图}

\subsection{点、边与度数}

图可以简单理解为：
\begin{itemize}
  \item 用\textbf{点}表示对象；
  \item 用\textbf{边}表示两个对象之间有某种关系。
\end{itemize}

例如：
\begin{itemize}
  \item 人与人认识，可以连一条边；
  \item 城市之间有公路，可以连一条边；
  \item 棋盘上的格点之间能一步走到，也可以连边。
\end{itemize}

一个点连着多少条边，这个数叫这个点的\textbf{度数}，记作 $\deg(v)$。

\begin{example}
在一个聚会中，5 个人两两握手。把每个人看成一个点，每次握手看成一条边。问一共有多少条边？
\end{example}

\begin{solution}
因为 5 个人两两都握手，所以任意两个人之间都有一条边。这就是一个 5 个点的完全图。

边数就是从 5 个人中任选 2 个人的方法数：
\[
\binom52=10.
\]

所以一共有
\[
10
\]
条边。
\end{solution}

\subsection{握手定理}

\begin{theorem}
在任意一个图中，所有点的度数之和等于边数的 2 倍，即
\[
\sum \deg(v)=2E,
\]
其中 $E$ 表示边数。
\end{theorem}

\begin{solution}
因为每一条边都会连接两个端点，所以当我们把所有点的度数加起来时，每条边都会被数两次。

因此总度数等于边数的 2 倍。
\end{solution}

\begin{skillbox}{握手定理特别适合处理这些题}
\begin{itemize}
  \item 已知每个点的度数，求边数；
  \item 已知边数，推度数和；
  \item 证明奇度点个数为偶数；
  \item 在“认识关系”“比赛安排”里快速建模。
\end{itemize}
\end{skillbox}

\begin{center}
\begin{tikzpicture}[
  scale=0.9,
  every node/.style={circle,draw=mainblue,fill=lightblue,inner sep=1.6pt,font=\small},
  edge/.style={draw=black!65,line width=0.9pt}
]
  \node[draw=none,fill=none,circle=false,font=\bfseries\color{mainblue}] at (2.0,3.0) {度数和就是把每条边的两个端点都数一遍};
  \node (A) at (0,1.2) {$2$};
  \node (B) at (1.5,2.4) {$3$};
  \node (C) at (3.0,1.4) {$2$};
  \node (D) at (2.4,-0.4) {$2$};
  \node (E) at (0.8,-0.6) {$1$};
  \foreach \u/\v in {A/B,B/C,C/D,D/A,B/D,A/E} {\draw[edge] (\u)--(\v);}
  \node[draw=none,fill=none,circle=false,font=\scriptsize] at (2.0,-1.5) {$2+3+2+2+1=10=2\times 5$};
\end{tikzpicture}
\end{center}

\section{用度数做计数}

\begin{example}
一个图有 8 个点，每个点的度数都等于 3。求这个图有多少条边。
\end{example}

\begin{solution}
总度数为
\[
8\times 3=24.
\]

由握手定理
\[
2E=24,
\]
所以
\[
E=12.
\]

因此这个图有
\[
12
\]
条边。
\end{solution}

\begin{example}
证明：任意一个图中，奇度数点的个数一定是偶数。
\end{example}

\begin{solution}
由握手定理，所有点度数之和等于
\[
2E,
\]
这是一个偶数。

如果一个和是偶数，那么其中奇数加数的个数必须是偶数；否则奇数个奇数相加会得到奇数。

因此图中奇度数点的个数一定是偶数。
\end{solution}

\section{完全图与二分图中的边数}

\subsection{完全图}

\begin{example}
完全图 $K_n$ 有多少条边？
\end{example}

\begin{solution}
完全图 $K_n$ 表示 $n$ 个点两两之间都有边。

所以边数就是从 $n$ 个点中任选 2 个点连成一条边的方法数：
\[
\binom{n}{2}=\frac{n(n-1)}{2}.
\]
\end{solution}

\begin{example}
9 个同学参加循环赛，每两人恰好比赛 1 场。一共要比多少场？
\end{example}

\begin{solution}
把每个同学看作一个点，每一场比赛看作一条边。这就是完全图 $K_9$ 的边数。

所以比赛场数为
\[
\binom92=36.
\]
\end{solution}

\subsection{二分图}

\begin{example}
有 4 个男生、5 个女生，如果规定只允许“男生和女生之间”连边，男生之间、女生之间都不连边，那么最多能连多少条边？
\end{example}

\begin{solution}
这是一个二分图。男生部分有 4 个点，女生部分有 5 个点。

若要边数最多，就让每个男生都和每个女生连边，于是边数为
\[
4\times 5=20.
\]
\end{solution}

\begin{theorem}
若一个二分图的两部分点数分别为 $a,b$，则它的边数最多为
\[
ab.
\]
\end{theorem}

\begin{solution}
因为边只能跨两部分连接，所以每个左边点最多与右边 $b$ 个点相连。总的最多边数就是
\[
ab.
\]
\end{solution}

\begin{center}
\begin{tikzpicture}[
  scale=0.88,
  every node/.style={circle,draw=mainblue,fill=lightblue,inner sep=1.5pt,font=\small},
  edge/.style={draw=black!60,line width=0.85pt}
]
  \node[draw=none,fill=none,circle=false,font=\bfseries\color{mainblue}] at (1.7,3.1) {完全图 $K_5$};
  \foreach \name/\x/\y in {A/0/1.3,B/1.0/2.5,C/2.5/2.0,D/2.7/0.5,E/1.1/-0.2} {
    \node (\name) at (\x,\y) {};
  }
  \foreach \u/\v in {A/B,A/C,A/D,A/E,B/C,B/D,B/E,C/D,C/E,D/E} {\draw[edge] (\u)--(\v);}
  \node[draw=none,fill=none,circle=false,font=\scriptsize] at (1.4,-1.0) {边数 $\binom52=10$};

  \begin{scope}[xshift=7.2cm]
    \node[draw=none,fill=none,circle=false,font=\bfseries\color{mainblue}] at (1.8,3.1) {二分图 $K_{4,5}$};
    \foreach \y/\name in {2.4/L1,1.6/L2,0.8/L3,0.0/L4} {\node (\name) at (0,\y) {};}
    \foreach \y/\name in {2.8/R1,2.0/R2,1.2/R3,0.4/R4,-0.4/R5} {\node (\name) at (3.6,\y) {};}
    \foreach \l in {L1,L2,L3,L4} {
      \foreach \r in {R1,R2,R3,R4,R5} {\draw[edge] (\l)--(\r);}
    }
    \node[draw=none,fill=none,circle=false,font=\scriptsize] at (1.8,-1.0) {边数最多 $4\times 5=20$};
  \end{scope}
\end{tikzpicture}
\end{center}

\section{网格图中的路径计数}

\begin{example}
在方格纸上，从左下角走到右上角，每次只能向右或向上走 1 格。若要从 $(0,0)$ 走到 $(3,2)$，共有多少条最短路径？
\end{example}

\begin{solution}
从 $(0,0)$ 到 $(3,2)$，一共要走 3 步向右、2 步向上，总共 5 步。

所以问题变成：在 5 个位置中选出哪 2 个位置放“向上”，或者选哪 3 个位置放“向右”。

因此最短路径条数为
\[
\binom52=10.
\]
\end{solution}

\begin{example}
若要从 $(0,0)$ 走到 $(4,4)$，每次只能向右或向上走 1 格，共有多少条最短路径？
\end{example}

\begin{solution}
要走 4 步向右和 4 步向上，总共 8 步。

所以最短路径条数为
\[
\binom84=70.
\]
\end{solution}

\begin{skillbox}{网格路径为什么也算图论计数}
网格上的每个格点可以看成一个点，相邻格点之间连边。于是“从起点到终点有多少条路”，本质上就是在图上数路径。
\end{skillbox}

\begin{center}
\begin{tikzpicture}[
  scale=0.9,
  every node/.style={circle,draw=black!70,fill=white,inner sep=1.2pt},
  edge/.style={draw=black!50,line width=0.8pt},
  path/.style={draw=accent,line width=1.3pt,-{Stealth[length=2.4mm]}}
]
  \node[draw=none,fill=none,circle=false,font=\bfseries\color{mainblue}] at (2.4,3.6) {网格最短路径 = 在若干步里安排“向右”和“向上”};
  \foreach \x in {0,1,2,3} {
    \foreach \y in {0,1,2} {\node (\x\y) at (1.2*\x,1.0*\y) {};}
  }
  \foreach \u/\v in {00/10,10/20,20/30,01/11,11/21,21/31,02/12,12/22,22/32,00/01,01/02,10/11,11/12,20/21,21/22,30/31,31/32} {
    \draw[edge] (\u)--(\v);
  }
  \draw[path] (00)--(10)--(11)--(21)--(22)--(32);
  \node[draw=none,fill=none,circle=false,font=\scriptsize] at (2.0,-0.8) {例：从 $(0,0)$ 到 $(3,2)$，共 3 步右、2 步上};
\end{tikzpicture}
\end{center}

\section{图论中的组合证明}

\begin{example}
某班有 7 位同学，每位同学都恰好认识 2 位同学。问“认识关系”一共对应多少对？
\end{example}

\begin{solution}
把每位同学看作一个点，“认识”看作一条无向边。

因为每位同学都认识 2 位同学，所以总度数是
\[
7\times 2=14.
\]

由握手定理
\[
2E=14,
\]
因此
\[
E=7.
\]

所以一共有
\[
7
\]
对认识关系。
\end{solution}

\begin{example}
证明：在任意 6 个人中，一定可以找到 3 个人，他们之间两两认识，或 3 个人，他们之间两两不认识。
\end{example}

\begin{solution}
这是一个非常经典的图论问题。把 6 个人看作 6 个点。

任选其中 1 个人 $A$，其余 5 个人与 $A$ 的关系只有两类：
\begin{itemize}
  \item 认识 $A$；
  \item 不认识 $A$。
\end{itemize}

由抽屉原理，这 5 个人中至少有 3 个人与 $A$ 的关系相同。

不妨设这 3 个人都认识 $A$，记为 $B,C,D$。

\begin{itemize}
  \item 如果 $B,C,D$ 中有两个人彼此认识，例如 $B$ 和 $C$ 认识，那么 $A,B,C$ 就两两认识；
  \item 如果 $B,C,D$ 中任意两个人都不认识，那么 $B,C,D$ 就构成 3 个互不认识的人。
\end{itemize}

两种情况总有一种发生，因此结论成立。
\end{solution}

\section{常见误区}

\begin{notebox}{图论计数里最常见的 5 个错误}
\begin{itemize}
  \item 把无向边数成了有向边，或者反过来；
  \item 用握手定理时忘记“每条边会被数两次”；
  \item 在完全图里把边数写成 $n(n-1)$，漏掉了除以 2；
  \item 在二分图里忘了边只能跨两部分连接；
  \item 网格路径题只会画图，不会转化成“若干步中选哪几步向上/向右”。
\end{itemize}
\end{notebox}

\section{课堂练习}

\begin{exercise}
完全图 $K_8$ 有多少条边？
\end{exercise}

\begin{exercise}
一个图有 10 个点，每个点的度数都等于 4。求这个图的边数。
\end{exercise}

\begin{exercise}
证明：图中不可能恰好有 3 个奇度数点。
\end{exercise}

\begin{exercise}
一个二分图两部分点数分别为 6 和 7，问边数最多是多少？
\end{exercise}

\begin{exercise}
从 $(0,0)$ 到 $(4,3)$，每次只能向右或向上走 1 格，共有多少条最短路径？
\end{exercise}

\begin{exercise}
8 个同学参加循环赛，每两人比赛 1 场，共有多少场比赛？
\end{exercise}

\newpage
\appendix
\section{练习解答}

\textbf{1}\begin{solution}
完全图 $K_8$ 的边数为
\[
\binom82=28.
\]
\end{solution}

\textbf{2}\begin{solution}
总度数为
\[
10\times 4=40.
\]

由握手定理
\[
2E=40,
\]
所以
\[
E=20.
\]
\end{solution}

\textbf{3}\begin{solution}
由握手定理，所有点的度数之和一定是偶数。

而奇数个奇数相加得到奇数，所以如果恰好有 3 个奇度数点，那么总度数和会是奇数，这与握手定理矛盾。

因此图中不可能恰好有 3 个奇度数点。
\end{solution}

\textbf{4}\begin{solution}
二分图两部分点数分别为 6 和 7 时，边数最多为
\[
6\times 7=42.
\]
\end{solution}

\textbf{5}\begin{solution}
从 $(0,0)$ 到 $(4,3)$，需要走 4 步向右、3 步向上，共 7 步。

所以最短路径数为
\[
\binom73=35.
\]
\end{solution}

\textbf{6}\begin{solution}
8 个同学两两比赛 1 场，就是完全图 $K_8$ 的边数，因此比赛场数为
\[
\binom82=28.
\]
\end{solution}

\end{document}
